\documentclass[12pt]{article}
\usepackage[utf8]{inputenc}
\usepackage[T1]{fontenc}
\usepackage{hyperref}
\title{Wilawila Language Resources for K-12 Curriculum Development}
\author{Compiled by Claude Sonnet 5.5 (claude-sonnet-5-5) for Andr\'as Kornai}
\date{}
\begin{document}
\maketitle

\section{Introduction}
Wilawila (Wila-Wila, Northern Language, Walmbi; ISO 639-3 wil) is a poorly documented Worrorran language of the northern Kimberley, Western Australia. It is probably no longer spoken. Wikipedia treats it as a dialect of Wunambal. The AIATSIS entry says the number of speakers is not known. Endangered Languages rates it critically endangered. It notes that people still identify with the language and are working to revitalize it.

\section{Primary Educational Resources}
None found. The AIATSIS entry lists no grammar.

\section{Linguistic and Reference Materials}
AIATSIS AUSTLANG entries: \url{https://collection.aiatsis.gov.au/austlang/language/K35} and \url{https://collection.aiatsis.gov.au/austlang/language/K44}. Endangered Languages profile: \url{https://endangeredlanguages.com/lang/3345}. Grambank: \url{https://grambank.clld.org/languages/wila1240}. Family context: \url{https://en.wikipedia.org/wiki/Worrorran_languages} and \url{https://en.wikipedia.org/wiki/Wunambal_language}. The search results cite Capell and Coate (1984) as the main source and a 1971 SIL word list. No URL was seen for either.

\section{Religious and Community Resources}
None found.

\section{Digital and Audio Resources}
None found.

\section{Curriculum Development Context}
A revitalization curriculum would rest on a very small documentary base. It would need to draw on archival word lists and on the closely related Wunambal materials, with guidance from descendants who identify with the language.
\end{document}